\begin{section}{Sheaves}
Sheaves are ubiquitous in mathematics and even though the reader might not had the opportunity
to apply them knowingly they can be sure to have used their structure implicitly in different categories.
From historic perspective, they originate in the domain of algebraic topology, where their ability to capture
consistenly the local behavior of objects such as continuous functions allowed an accessable
treatment of rather involved topological spaces. Over the course of time one of the greatest 
extensions of this \textit{tracking idea} has been achieved by Alexander Grothendieck, who used
them in his work EGA \cite{egatransl} in their specialized version as schemes to generalize algebraic geometry to the spectrum 
of rings, unifying geometric perspectives from the polynomial and number theoretic settings.
The sheaves in this work have a similar spirit, so without further ado:

\begin{definition}[Presheaves and Sheaves]
    Given a topological space \( X \), a sheaf is a structure s.t. 
    for all open \( U \subseteq X \), there is a map of sets \( \mathcal{O} : U \to \mathcal{O}(U) \) 
    and for all open \( V \subseteq U \) with their respective inclusion \( \iota : V \to U  \),
    corresponds a map \( r_{U,V} : \mathcal{O}(U) \to \mathcal{O}(V) \).
    The \( r \) stands for restriction and we use the notation \( f|_V := r_{U, V}(f) \) for \( f \in \mathcal{O}(U) \). Additionally 
    the maps should satisfy:
    \begin{enumerate}
        \item \( r_{U, U} = \text{id}_{\mathcal{O}(U)} \)
        \item If \( U \subseteq W \) then \( r_{U, V} \circ r_{W, U} = r_{W, V} \).
        \item (Locality): For a cover \( \bigcup_i U_i = U \) and  \( f, g \in \mathcal{O}(U) \), s.t. \( f|_{U_i} = g|_{U_i} \) 
            for all \( i \), then \( f = g \).
        \item (Gluing): For a cover \( \bigcup_i U_i = U \) and a collection \( \{ f_i \}_i,\ f_i \in \mathcal{O}(U_i) \), 
           s.t. \( f_i|_{U_i \cap U_j} = f_j|_{U_i \cap U_j} \), then there exists \( f \in \mathcal{O}(U) \)
           s.t. \( f|_{U_i} = f_i \).
    \end{enumerate}
    A structure which only satisfies the first two properties is a presheaf. If it also satisfies locality it is separated.
\end{definition} 
%The quite general definition hints already at the fact that sheaves are ubiquitous in mathematics, e.g. smooth functions over a manifold
%form a sheaf, the sheaf of Schwartz distributions over a space and there is even a sheaf of measures on Borel spaces \cite{sheafmeasures}.
%In the affine case here, there is a particular choice of sheaf, which got its own name. For this we also need to define morphisms.

\begin{example} \ref{algstructuresheaf}
    Consider the affine space \( \aff{n}{\mathbb{K}} = \mathbb{K}^n\) over some field \( \mathbb{K} \) 
    and the associated commutative ring of multivariate polynomials \( \mathbb{K}[\mathbf{X}] \).
    One can define the Zariski topology (a \( T_1 \)-topology) over \( \aff{n}{\mathbb{K}} \), whose subbasis of closed sets is induced through 
    \( J \subseteq \mathbb{K}[\mathbf{X}] \) via their vanishing sets, that is 
    \[ V(J) = \{\, a \in \mathbb{K}^n : f(a) = 0 \text{ for all } f \in J \,\}. \]
    Clearly, one also obtains a subbasis of open sets by taking the complements of vanishing sets (denoted by \( D(J) \)
    traditionally), which are points over which collection of polynomials do not vanish.
    Define then for an open set \( U \subset \aff{n}{\mathbb{K}} \) the set of \( \mathcal{O}(U) \) regular functions over \( U \), that is 
    \( f : U \to \mathbb{K} \), s.t. for every \( x \in U \) there is some open neighborhood \( V \subset U \) of it on which it holds
    \( f|_{V} \neq 0 \) and \( f(x) = \frac{h(x)}{g(x)} \) for \( h, g \in \mathbb{K}[\mathbf{X}] \).
    This mapping \( \mathcal{O} : \text{Open}(\aff{n}{\mathbb{K}}) \to \mathcal{O}(U) \) defines in fact a sheaf over \( \aff{n}{\mathbb{K}} \)
    with the restriction being the standard domain restriction of functions as the reader can verify. Moreover, the 
    \( \mathcal{O}(U) \) are also rings, since adding or multiplying two non-vanishing functions \( f_1, f_2 \) yields another one 
    over the intersection \( U_1 \cap U_2 \) of their domains, turning \( \aff{n}{\mathbb{K}} \) into a locally ringed space.
\end{example}

This example hopefully illuminates how the construction of the sheaf of smooth functions over a smooth manifold would look like.
As is apparent from the Definition, the local behavior of functions can be tracked through the restriction morphism. If we continue this 
zoom into smaller and smaller open sets indefinitely and end up at a point (which is more than often not an open set in the topology of the underlying space \( X \)) the result is often something similar to an evaluation. The formal construction of this idea requires the colimit 
from category theory, which can be found in the Appendix \ref{defcolimit}, the main challenge is to properly handle that it doesn't matter
from which open set we initialize our restrictions to the point. This idea has its own name for sheaves: 

\begin{definition}
    The stalk at a point \( x \in X \) is the colimit of the sheaf sections \( \mathcal{O}(U) \) of the open neighborhoods \( U \) of the point,
    denoted by
    \[
        \mathcal{O}_{X,x} = \varinjlim_{U\ni x} \mathcal{F}(U).
    \]
\end{definition}
\begin{example}
    For the sheaf of regular functions over \( \aff{n}{\mathbb{K}} \) the stalk at a point \( x \) is given by the ring 
    of rational polynomial functions \( f = \frac{g}{h} \) with \( h(x) \neq 0 \).
\end{example}
\end{section}

Everything is a matter of perspective and there is rarely one canonical choice over some space and its 
topology. Already for smooth manifolds we can assign for each \( k \in \mathbb{N} \) to an open set
the set of \( k \)-differentiable functions over it, giving a plethora of sheaves. It can be useful to inspect
how this different sheaves relate to each other as long as there is a consistent way of translating between their sections
which is made precise by the following idea.
\begin{definition}[Morphism of sheaves] \label{sheafmorphism}
    A morphism \(\varphi : \mathcal{F} \to \mathcal{G}\) of sheaves on \(X\) is a rule which assigns to each open \(U \subseteq X\) 
    a map of sets \(\varphi : \mathcal{F}(U) \to \mathcal{G}(U)\) compatible with restriction maps, i.e., whenever \(V \subseteq U \subseteq X\) are open 
    the following diagram commutes: 
    \[
        \begin{tikzcd}
            \mathcal{F}(U) \arrow[r, "\varphi"] \arrow[d, "r_{U,V}"] & \mathcal{G}(U) \arrow[d, "r_{U,V}"] \\
            \mathcal{F}(V) \arrow[r, "\varphi"] & \mathcal{G}(V)
        \end{tikzcd}
    \]
    An isomorphism of sheaves is a morphism \( \varphi : \mathcal{F} \to \mathcal{G} \), s.t there exists a morphism 
    \( \varphi^{-1} : \mathcal{G} \to \mathcal{F} \), satisfying 
    \begin{align*}
        \varphi^{-1} \circ \varphi &= \text{id}_{(X, \mathcal{F})} \\
    \varphi \circ \varphi^{-1} &= \text{id}_{(Y, \mathcal{G})} 
    \end{align*}
\end{definition}
Note that if there exists a homeomorphism \( \pi : X \to Y \) between topological spaces, we can
also define morphisms of sheaves with different underlying spaces, namely if \( (X, \mathcal{F}) \) and \( (Y, \mathcal{G}) \)
are sheaves, then for a open sets \( U \subseteq X, V \subseteq Y \) one could construct maps 
\( \varphi : \mathcal{F}(U) \to \mathcal{G}(\pi(U)) \), \( \varphi' : \mathcal{F}(\pi^{-1}(V)) \to \mathcal{G}(V)  \)
and if for all \( U \) (respectively \( V \)) these maps satisfy the definition of \ref{sheafmorphism},
we get a morphism between sheaves with different underlying space. Unfortunately, these can't be the same
(simply because \( X \) and \( Y \) are not the same space), but there's a one to one correspondence
(due the adjointness of the preimage sheaf and pushfoward sheaf functor). In case that \( \pi \)
is just a continuous map, one refers typically to morphisms of sheaves to the direction \( \varphi' : \mathcal{F}(\pi^{-1}(V)) \to \mathcal{G}(V) \) 
(so it is a morphism of sheaves on \( Y \)).
